\chapter{The templates}
\label{app:templates}

Ten fill-in documents live in \code{templates/}. Copy the directory per
company; the originals stay clean.

Each carries an \code{UNFILLED} marker at the top. A build gate reports how
many templates still carry it, which is not a criticism \dash{} it is so that a
half-filled document is visible rather than mistaken for a finished one.

\begin{kittable}{@{}llX@{}}
\textbf{File} & \textbf{Chapter} & \textbf{Fill it in} \\
\midrule
\code{fit-check.md} & \ref{ch:fit} & before applying. Twenty minutes \\
\code{company-brief.md} & \ref{ch:company} & before the screen \\
\code{role-rubric.md} & \ref{ch:segment} & before the screen \\
\code{interviewer-brief.md} & \ref{ch:person} & one per person, before each call \\
\code{claim-inventory.md} & \ref{ch:claims} & before you submit anything \\
\code{banned-claims.md} & \ref{ch:claims} & before every conversation \\
\code{project-deep-dive.md} & \ref{ch:deep-dive} & once, then reuse \\
\code{run-sheet.md} & \ref{ch:runsheet} & before a build round \\
\code{decision-log.md} & \ref{ch:runsheet} & during a build round \\
\code{retrospective.md} & \ref{ch:after} & after every process \\
\code{outcome-ledger.md} & \ref{ch:after} & one row per process, forever \\
\end{kittable}

\section{The order}

Three come before you apply: \code{fit-check}, \code{role-rubric},
\code{company-brief}. In that order, because the first can end the exercise.

Two come before you submit: \code{claim-inventory}, and the artefact's own
deviations register from Chapter~\ref{ch:ledger}.

Two come before every conversation: \code{interviewer-brief} and
\code{banned-claims}. The second takes five minutes after the first time.

Two come with a build round. Fill in \code{run-sheet.md} with absolute times
beforehand, and keep \code{decision-log.md} open during.

Two come after: \code{retrospective}, and a row in \code{outcome-ledger}.

\section{The one that compounds}

\code{outcome-ledger.md} is the only document here that is worth more the
longer you keep it. One row per process: stages reached, where it ended, and
\dash{} the column that matters \dash{} which of Chapter~\ref{ch:classes}'s
evidence classes it ended in.

\judgement{A single process tells you almost nothing. Outcomes at this level
are noisy and one rejection is weak evidence about anything. Four rows with a
pattern in the last column tell you where your actual constraint is, and
nothing else will.}
